\documentclass{exam}
\usepackage{../../shah350}

\hypersetup{colorlinks=true, linktoc=section, linkcolor=blue}

\pagestyle{headandfoot}
\firstpageheadrule
\runningheadrule
\firstpageheader{Prof. Lepowsky \\ Linear Algebra}{Homework \#3}{Jeevan Shah}
\runningheader{Linear Algebra \\ Homework \#3}{}{Shah}
\firstpagefooter{}{\thepage}{}
\runningfooter{ }{\thepage}{ }

\printanswers


\begin{document}
  \colorbox{red}{\underline{$1.6.2b$}:}
  \begin{solution}
    Since there are $3$ vectors in the set and $\dim \RR^3 = 3$ it suffices to check for linear independence by Corollary 2:
    \begin{align*}
      A = \begin{pmatrix}
        2 & 0 & 6 \\
        -4 & 3 & 0 \\
        1 & -1 & -1
      \end{pmatrix}
      \xrightarrow{2_1 \to \frac{1}{2}r_1}
      &\begin{pmatrix}
        1 & 0 & 3 \\
        -4 & 3 & 0 \\
        1 & -1 & 1 
      \end{pmatrix} \\
      \xrightarrow[r_3 \to r_3 - r_1]{r_2 \to r_2 + 4r_1}
      &\begin{pmatrix}
        1 & 0 & 3 \\
        0 & 3 & 12 \\
        0 & -1 & -4  
      \end{pmatrix} \\
      \xrightarrow{r_3 \to 3r_3 + r_2} 
      &\begin{pmatrix}
        1 & 0 & 3 \\
        0 & 3 & 12 \\
        0 & 0 & 0
      \end{pmatrix}
    \end{align*} 
    thus $\rank A = 2 < 3$ and hence the set is linearly dependent and thus not a base for $\RR^3$.
  \end{solution}

  \colorbox{red}{\underline{$1.6.3b$}:}
  \begin{solution}
    Using the fact that $P_2(\RR) \cong \RR^3$ to write the vectors as column vectors in $\RR^3$ and then check for linear independence this way:
    \begin{align*}
      A = \begin{pmatrix}
        1 & 3 & 0 \\
        2 & 0 & 1 \\
        1 & 1 & 1
      \end{pmatrix}
      \xrightarrow[r_3 \to r_3 - r_1]{r_2 \to r_2 - 2r_1}
      &\begin{pmatrix}
        1 & 3 & 0 \\
        0 & -6 & 1 \\
        0 & -2 & 1 
      \end{pmatrix}
      \xrightarrow{r_3 \to 3r_3 - r_2}
      \begin{pmatrix}
        1 & 3 & 0\\
        0 & -6 & 1 \\
        0 & 0 & 2
      \end{pmatrix}
    \end{align*}
    thus $\rank A = 3$ and so the set is linearly independent and thus a basis by Corollary 2.
  \end{solution}

  \colorbox{red}{\underline{$1.6.9$}:}
  \begin{solution}
    \begin{align*}
      \begin{amatrix}{4}
        1 & 0 & 0 & 0 & a_1 \\
        1 & 1 & 0 & 0 & a_2 \\
        1 & 1 & 1 & 0 & a_3 \\
        1 & 1 & 1 & 1 & a_4
      \end{amatrix}
      \xrightarrow[r_4 \to r_4 - r_1]{\substack{r_2 \to r_2 - r_1 \\ r_3 \to r_3 - r_1}}
      &\begin{amatrix}{4}
        1 & 0 & 0 & 0 & a_1 \\
        0 & 1 & 0 & 0 & a_2 - a_1 \\
        0 & 1 & 1 & 0 & a_3 - a_1 \\
        0 & 1 & 1 & 1 & a_4 - a_1
      \end{amatrix} \\
      \xrightarrow[r_4 \to r_4 - r_2]{r_3 \to r_3 - r_2}
      &\begin{amatrix}{4}
        1 & 0 & 0 & 0 & a_1 \\
        0 & 1 & 0 & 0 & a_2 - a_1 \\
        0 & 0 & 1 & 0 & a_3 - a_2 \\
        0 & 0 & 1 & 1 & a_4 - a_2
      \end{amatrix} \\
      \xrightarrow{r_4 \to r_4 - r_3}
      &\begin{amatrix}{4}
        1 & 0 & 0 & 0 & a_1 \\
        0 & 1 & 0 & 0 & a_2 - a_1 \\
        0 & 0 & 1 & 0 & a_3 - a_2 \\
        0 & 0 & 0 & 1 & a_4 - a_3
      \end{amatrix} \\
    \end{align*}
    Thus,  
    \[
      \begin{pmatrix}
        a_1 \\ a_2 \\ a_3 \\ a_4
      \end{pmatrix} 
      = a_1\vec{u_1} + (a_2 - a_1)\vec{u_2} + (a_3 - a_2)\vec{u_3} + (a_4 - a_3)\vec{u_4}
    \]
  \end{solution}

  \colorbox{red}{\underline{$2.1.2$}:}
  \begin{solution}
    $T$ is of the form $L_A$ for 
    \[
      A = \begin{pmatrix}
        1 & -1 & 0 \\
        0 & 0 & 2
      \end{pmatrix}.
    \]
    To show $T$ is linear, first consider $\vec{u} = (a_1, a_2, a_3)$ and $\vec{v} = (b_1, b_2, b_3)$. Then
    \begin{align*}
      &T(\vec{u} + \vec{v}) = \begin{pmatrix}
        (a_1 + b_1) - (a_2 + b_2) \\
        2(a_3 + b_3)
      \end{pmatrix}
      = \begin{pmatrix}
        a_1 - a_2 \\
        2a_3
      \end{pmatrix}
      + \begin{pmatrix}
        b_1 - b_2 \\
        2b_3
      \end{pmatrix}
      = T(\vec{u}) + T(\vec{v}) \\
      &T(c\vec{u}) = \begin{pmatrix}
        c(a_1 - a_2) \\ 2ca_3
      \end{pmatrix}
      = c\begin{pmatrix}
        a_1 - a_2 \\ 2a_3
      \end{pmatrix} 
      = cT(\vec{u})
    \end{align*}
    and so $T$ is linear. To find a basis for $N(T)$ we see 
    \[
      A = \begin{pmatrix}
        1 & -1 & 0 \\
        0 & 0 & 2
      \end{pmatrix} 
      \Rightarrow \begin{cases}
        x_1 - x_2 = 0 \Rightarrow x_1 = x_2 \\
        2x_3 = 0 \Rightarrow x_3 = 0
      \end{cases}
      \Rightarrow X = \begin{pmatrix}
        x_1 \\ x_1 \\ 0
      \end{pmatrix}
      \Rightarrow \Span{\begin{pmatrix}
        1 \\ 1 \\ 0
      \end{pmatrix}} = N(T)
    \]
    and so $\braces{\begin{pmatrix} 1 \\ 1 \\ 0 \end{pmatrix}}$ is a basis for $N(T)$. For $R(T)$ we have that
    \[
      \braces{\begin{pmatrix}
        1 \\ 0
      \end{pmatrix}, \begin{pmatrix}
        0 \\ 2
      \end{pmatrix}}
    \]
    is a basis for $R(T)$ since they are the vectors corresponding to the pivot columns of RREF$(A)$. Thus 
    \[
      \nul{T} + \rank{T} = 1 + 2 = 3 = \dim \RR^3
    \]
    verifying the dimension theorem. Since $N(T) \neq \braces{0}$, $T$ is not one-to-one by Thm2.4. $T$ is onto since $\rank{A} = 2 = \dim \RR^2$.
  \end{solution}

  \colorbox{red}{\underline{$2.1.3$}:}
  \begin{solution}
    T is of the form $L_A$ for 
    \[
      A = \begin{pmatrix}
        1 & 1 \\
        0 & 0 \\
        2 & -1 
      \end{pmatrix}.
    \]
    To see that $T$ is linear consider $\vec{u}, \vec{v} \in \RR^2$ such that $\vec{u} = (a_1, a_2), \vec{v} = (b_1, b_2)$, Then, 
    \begin{align*}
      &T(\vec{u} + \vec{v}) = \begin{pmatrix}
        a_1 + b_1 + a_2 + b_2 \\
        0 \\
        2(a_1 + b_1) - (a_2 + b_2)
      \end{pmatrix}
      = \begin{pmatrix}
        a_1 + a_2 \\
        0 \\
        2a_1 - a_2
      \end{pmatrix}
      + \begin{pmatrix}
        b_1 + b_2 \\
        0 \\
        2(b_1 - b_2)
      \end{pmatrix}
      = T(\vec{u}) + T(\vec{v}) \\
      &T(c\vec{u}) = \begin{pmatrix}
        c(a_1 + a_2) \\
        0 \\
        2ca_1 - ca_2
      \end{pmatrix}
      = c\begin{pmatrix}
        a_1 + a_2 \\
        0 \\
        2a_1 - a_2
      \end{pmatrix}
      = cT(\vec{u})
    \end{align*}
    and so $T$ is linear. To find a basis for $N(T)$ we have 
    \[
      A = \begin{pmatrix}
        1 & 1 \\
        0 & 0 \\
        2 & -1 
      \end{pmatrix} 
      \xrightarrow{r_2 \to r_2 - 2r_1}
      \begin{pmatrix}
        1 & 1 \\
        0 & 0 \\
        0 & -3 
      \end{pmatrix}
      \Rightarrow \begin{cases}
        -3x_2 = 0 \Rightarrow x_2 = 0 \\
        x_1 + x_2 = 0 \Rightarrow x_1 = -x_2
      \end{cases}
      \Rightarrow X = \begin{pmatrix}
        0 \\ 0
      \end{pmatrix}
    \]
    and hence $N(T) = \braces{0}$. For $R(T)$ we know from above that the first and second columns are pivot columns and so 
    \[
      \braces{
        \begin{pmatrix}
          1 \\ 0 \\ 2
        \end{pmatrix},
        \begin{pmatrix}
          1 \\ 0 \\ -1
        \end{pmatrix}
      }
    \]
    is a basis for $R(T)$. Since $N(T) = \braces{0}$, $T$ is one-to-one and because $\rank{A} = 2 < 3 = \dim \RR^3$ T is not onto.
  \end{solution}

  \colorbox{red}{\underline{$2.1.9$}:}
  \begin{solution}
    \begin{parts}
      \part $T$ is not linear since for $\vec{u} = (a_1, a_2)$ and $\vec{v} = (b_1, b_2)$ we have 
      \[
        T(\vec{u} + \vec{v}) = \begin{pmatrix}
          1 \\ a_2 + b_2
        \end{pmatrix}
        \neq \begin{pmatrix}
          2 \\ a_2 + b_2
        \end{pmatrix}
        = T(\vec{u}) + T(\vec{v})
      \]

      \part $T$ is not linear since for $\vec{u} = (a_1, a_2)$ and $c \in \RR$ we have 
      \[
        T(c\vec{u}) = \begin{pmatrix}
          ca_1 \\ (ca_2)^2
        \end{pmatrix}
        \neq 
        \begin{pmatrix}
          ca_1 \\
          c(a_1)^2
        \end{pmatrix}
        = cT(\vec{u})
      \]
    \end{parts}
  \end{solution}

  \colorbox{red}{\underline{$2.1.15$}:}
  \begin{solution}
    \begin{proof}
      We first show that $T$ is linear. Consider $f, g \in P(\RR)$ and $c \in \RR$. Then, 
      \begin{align*}
        &T(f(x) + g(x)) = \int_{0}^{x} f(t)\, \mathrm{d}t + \int_{0}^{x} g(t)\, \mathrm{d}t = \int_{0}^{x} f(t) + g(t)\, \mathrm{d}t = T(f(x) + g(x)) \\
        &T(cf(x)) = \int_{0}^{x} cf(t)\,\mathrm{d}t = c\int{0}^{x} f(t)\,\mathrm{d}t = cT(f(x))
      \end{align*}
      and hence $T$ is linear. It is clear that $T$ is not onto as any constant polynomial ($f(x) = n$ for $n \in \RR$) is not in $R(T)$ since there is nothing that integrates to just a constant. We show that $T$ is one-to-one by assuming that $T(f(x)) = T(g(x))$ and showing $f = g$: 
      \begin{align*}
        T(f(x)) = T(g(x)) &\Rightarrow \int_{0}^{x} f(t)\, \mathrm{d}t = \int_{0}^{x} g(t)\,\mathrm{d}t \\
        &\Rightarrow \int_{0}^{x} f(t) - g(t)\,\mathrm{d}t = 0 \\
        &\Rightarrow \frac{\mathrm{d}}{\mathrm{d}x}\left[\int_{0}^{x}f(t) - g(t)\,\mathrm{d}t \right] = 0 \\
        &\Rightarrow f(x) - g(x) = 0 \\
        &\Rightarrow f(x) = g(x) \quad \forall\, x \in \RR
      \end{align*}
      as was to be shown. 
    \end{proof}
  \end{solution}

  \newpage

  \colorbox{red}{\underline{$2.2.2$}:}
  \begin{solution}
    \begin{parts}
      \part We have that 
      \[
        T(e_1) = \begin{pmatrix}
          2 \\ 3 \\ 1
        \end{pmatrix}
        \quad\text{and}\quad
        T(e_2) = \begin{pmatrix}
          -1 \\ 4 \\ 0
        \end{pmatrix}
      \]
      so 
      \[
        \left[T\right]_{\beta}^{\gamma} = \begin{pmatrix}
          2 & 1 \\
          3 & 4 \\
          1 & 0 
        \end{pmatrix} = A
      \]
      and so the matrix representation fo $T$ in the standard ordered bases for $\RR^{2}$ and $\RR^{3}$ is the same as the matrix presentation of $T$ in the form $T=L_A$.

      \part We have that 
      \[
        T(e_1) = \begin{pmatrix}
          2 \\ 1
        \end{pmatrix}
        \quad\text{and}\quad
        T(e_2) = \begin{pmatrix}
          3 \\ 0
        \end{pmatrix}
        \quad\text{and}\quad
        T(e_3) = \begin{pmatrix}
          -1 \\ 1
        \end{pmatrix}
      \]
      so 
      \[
        \left[T\right]_{\beta}^{\gamma} = \begin{pmatrix}
          2 & 3 & -1 \\
          1 & 0 & 1
        \end{pmatrix} = A
      \]
      and so the matrix representation fo $T$ in the standard ordered bases for $\RR^{2}$ and $\RR^{3}$ is the same as the matrix presentation of $T$ in the form $T=L_A$.

      \part We have that 
      \[
        T(e_1) = 2
        \quad\text{and}\quad
        T(e_2) = 1
        \quad\text{and}\quad
        T(e_3) = -3
      \]
      so 
      \[
        \left[T\right]_{\beta}^{\gamma} = \begin{pmatrix}
          2 & 1 & -3
        \end{pmatrix} = A
      \]
      and so the matrix representation fo $T$ in the standard ordered bases for $\RR$ and $\RR^{3}$ is the same as the matrix presentation of $T$ in the form $T=L_A$.
    \end{parts}
  \end{solution}


  \colorbox{red}{\underline{$2.2.4$}:}
  \begin{solution}
    We have that 
    \[
      T(e_1) = 1 
      \quad\text{and}\quad
      T(e_2) = 1+x^2
      \quad\text{and}\quad 
      T(e_3) = 0
      \quad\text{and}\quad 
      T(e_4) = x
    \]
    so 
    \[
      \left[T\right]_{\beta}^{\gamma} = \begin{pmatrix}
        1 & 1 & 0 & 0 \\
        0 & 0 & 0 & 1 \\
        0 & 1 & 0 & 0 
      \end{pmatrix}
    \]
  \end{solution}

  \colorbox{red}{\underline{$2.2.5d$}:}
  \begin{solution}
    We have that 
    \[
      T(e_1) = 1 
      \quad\text{and}\quad 
      T(e_2) = 2 
      \quad\text{and}\quad
      T(e_3) = 4  
    \]
    so 
    \(
      \left[T\right]_{\beta}^{\gamma} = \begin{pmatrix}
        1 & 2 & 4
      \end{pmatrix}
    \)
  \end{solution}
\end{document}